% GENERATED FILE. Edit canonical lessons, then run npm run generate:books.
% canonical-source-hash: fnv1a64:8ad659c9ef769b16
% canonical-lessons: TA-C68-shop, TA-C68-price, TA-W24-read-kadai, TA-C68-money, TA-C68-buy, TA-C68-bag

\chapter{At the Shop}
\label{ch:ta-at-the-shop}
\hlchaptermodality{68}

\begin{chapteropening}
\textbf{By the end of this chapter:} \emph{I can name a shop, a price, money, buying and a bag in Tamil, and ask what something costs.}

Complete TA-C68-bag and say all five words of the chapter in order.
\end{chapteropening}

\section[kaḍai]{\ta{கடை} (kaḍai) — a shop}

\label{lesson:TA-C68-shop}



\begin{quote}
Before the new one: what did \emph{toppi} mean?
\end{quote}

\subsection*{You'll want to know: \ta{கடை}}
\textbf{\ta{கடை}} (\emph{kaḍai}) — \textquotedblleft{}a shop\textquotedblright{}.

Any shop at all: the tea shop, the cloth shop, the one at the end of the street. Tamil does not split the word by what is sold — a \ta{கடை} is a \ta{கடை}, and what is inside it is said separately.

It ends in the same \textbf{\ta{ை}} as \ta{தலை} and \ta{சட்டை}, and it is where every garment of the last chapter came from.

\begin{grammarlens}[title={What you've built}]
The place the next four words happen in.
\end{grammarlens}

\subsection*{Your turn}
\begin{itemize}
\raggedright
  \item \emph{Say it:} \emph{kaḍai}
  \item \emph{Say it:} \emph{kaḍai}, then name one thing from last chapter you would go there for
  \item \emph{Recall:} say \emph{seruppu}, then say \emph{toppi}, then say \emph{kaḍai}
\end{itemize}

\begin{tcolorbox}[breakable,colback=teal!4,colframe=teal!35!black,title={Before you move on}]
Say it once more.
\end{tcolorbox}
\section[vilai]{\ta{விலை} (vilai) — a price}

\label{lesson:TA-C68-price}



\begin{quote}
Before the new one: what did \emph{kaḍai} mean?
\end{quote}

\subsection*{You'll want to know: \ta{விலை}}
\textbf{\ta{விலை}} (\emph{vilai}) — \textquotedblleft{}a price\textquotedblright{}.

The number the shopkeeper says back to you.

This one word turns a question you already own into a whole sentence. You have \ta{எவ்வளவு} (\emph{evvaḷavu}), \textquotedblleft{}how much\textquotedblright{}. Put them together and you have \textbf{\ta{விலை} \ta{எவ்வளவு}?} — \textquotedblleft{}the price, how much?\textquotedblright{}

Notice what is not there: no verb, and no word for \textquotedblleft{}is\textquotedblright{}. Tamil names the thing, then asks about it, and stops.

\begin{grammarlens}[title={What you've built}]
Two words, and a question you can already ask.
\end{grammarlens}

\subsection*{Your turn}
\begin{itemize}
\raggedright
  \item \emph{Say it:} \emph{vilai}
  \item \emph{Say it:} \emph{vilai evvaḷavu?}
  \item \emph{Say it:} \emph{kaḍai}, then \emph{vilai}, then the whole question
  \item \emph{Recall:} say \emph{toppi}, then say \emph{kaḍai}, then say \emph{vilai}
\end{itemize}

\begin{tcolorbox}[breakable,colback=teal!4,colframe=teal!35!black,title={Before you move on}]
How do you ask what something costs? (\textbf{\ta{விலை} \ta{எவ்வளவு}?})
\end{tcolorbox}
\section[kaḍai]{\ta{கடை} — a sign that wraps the letter}

\label{lesson:TA-W24-read-kadai}



\begin{quote}
Before the page: you can already say \emph{kaḍai} and \emph{vilai}. Now read the first one.
\end{quote}

\subsection*{Script you'll notice: \ta{கடை} — a sign that wraps the letter}
\textbf{\ta{கடை}} is as short as a Tamil word gets: two letters, one sign, nothing new.

\noindent
\begin{tabularx}{\linewidth}{@{}>{\raggedright\arraybackslash}X>{\raggedright\arraybackslash}X>{\raggedright\arraybackslash}X@{}}
\toprule
\textbf{piece} & \textbf{} & \textbf{} \\
\midrule
\textbf{\ta{க}} & \emph{ka} & the letter with its own built-in \emph{a} \\
\textbf{\ta{டை}} & \emph{ṭai} & \ta{ட} with the \textbf{\ta{ை}} sign, the same one you read at the end of \ta{சட்டை} \\
\bottomrule
\end{tabularx}

\begin{quote}
\textbf{\ta{க} · \ta{டை} $\to$ \ta{கடை}}
\end{quote}

Look at where the \textbf{\ta{ை}} sits. It is drawn to the \textbf{left} of its letter but said \textbf{after} it, which is the habit you first met on the \ta{ே} of \ta{மேலே}. Two signs now, both hanging in front of the letter they belong to.

And notice what the same two letters do without the sign: \textbf{\ta{கடு}}, \textbf{\ta{கடி}} and \textbf{\ta{கடை}} are three different words with the same skeleton. In Tamil the sign is not decoration.

\subsection*{Your turn}
\begin{itemize}
\raggedright
  \item \emph{Read it:} \textbf{\ta{க}} — then \textbf{\ta{டை}} — then \textbf{\ta{கடை}}
  \item \emph{Point to:} the \textbf{\ta{ை}}, and say which side of its letter it is drawn on
  \item \emph{Read it:} \textbf{\ta{சட்டை}}, then \textbf{\ta{கடை}}, and find the ending they share
  \item \emph{Recall:} say \emph{kaḍai}, then say \emph{vilai}, then read \textbf{\ta{கடை}}
\end{itemize}

\begin{tcolorbox}[breakable,colback=teal!4,colframe=teal!35!black,title={Before you move on}]
Read \textbf{\ta{கடை}}. On which side of \ta{ட} is the \textbf{\ta{ை}} drawn, and when is it said? (Drawn on the \textbf{left}; said \textbf{after}.)
\end{tcolorbox}
\section[paṇam]{\ta{பணம்} (paṇam) — money}

\label{lesson:TA-C68-money}



\begin{quote}
Before the new one: read \textbf{\ta{கடை}}, and say what it means.
\end{quote}

\subsection*{You'll want to know: \ta{பணம்}}
\textbf{\ta{பணம்}} (\emph{paṇam}) — \textquotedblleft{}money\textquotedblright{}.

Money, in the plain everyday sense — what is in your hand, not what is in a bank.

The \ta{ண} in the middle is the curled-back one, the same \ta{ண} you have been saying since \ta{வணக்கம்} on the very first page. And it closes with \textbf{\ta{ம்}}, the ending that finishes so many Tamil nouns: \ta{வணக்கம்}, \ta{தூக்கம்}, and now \ta{பணம்}.

\begin{grammarlens}[title={What you've built}]
Three: the shop, the price, and what you pay it with.
\end{grammarlens}

\subsection*{Your turn}
\begin{itemize}
\raggedright
  \item \emph{Say it:} \emph{paṇam}
  \item \emph{Say it:} \emph{vaṇakkam}, then \emph{paṇam} — and say what the two have in common
  \item \emph{Recall:} say \emph{vilai}, then read \textbf{\ta{கடை}}, then say \emph{paṇam}
\end{itemize}

\begin{tcolorbox}[breakable,colback=teal!4,colframe=teal!35!black,title={Before you move on}]
What does \ta{பணம்} mean? (\textquotedblleft{}Money\textquotedblright{}.)
\end{tcolorbox}
\section[vāṅgu]{\ta{வாங்கு} (vāṅgu) — to buy, to get}

\label{lesson:TA-C68-buy}



\begin{quote}
Before the new one: what did \emph{paṇam} mean?
\end{quote}

\subsection*{You'll want to know: \ta{வாங்கு}}
\textbf{\ta{வாங்கு}} (\emph{vāṅgu}) — \textquotedblleft{}to buy, to get\textquotedblright{}.

The verb that makes a shop work.

It is also the ordinary Tamil verb for \textbf{receiving}. Buying something and being handed something are the same act here, and one word covers both — so \ta{வாங்கு} is as useful at a doorway as at a counter.

Set it against \ta{கொடு} (\emph{koṭu}), \textquotedblleft{}give\textquotedblright{}, which you already have. One hand out, one hand in, and the pair covers most of what happens over a counter.

\begin{grammarlens}[title={What you've built}]
Four, and now something can change hands.
\end{grammarlens}

\subsection*{Your turn}
\begin{itemize}
\raggedright
  \item \emph{Say it:} \emph{vāṅgu}
  \item \emph{Say it:} \emph{koṭu}, then \emph{vāṅgu} — and say which way each hand is moving
  \item \emph{Recall:} read \textbf{\ta{கடை}}, then say \emph{paṇam}, then say \emph{vāṅgu}
\end{itemize}

\begin{tcolorbox}[breakable,colback=teal!4,colframe=teal!35!black,title={Before you move on}]
What two jobs does \ta{வாங்கு} do? (\textbf{Buying}, and \textbf{receiving}.)
\end{tcolorbox}
\section[pai]{\ta{பை} (pai) — a bag}

\label{lesson:TA-C68-bag}



\begin{quote}
Before the new one: what did \emph{vāṅgu} mean, and what did \emph{paṇam} mean?
\end{quote}

\subsection*{You'll want to know: \ta{பை}}
\textbf{\ta{பை}} (\emph{pai}) — \textquotedblleft{}a bag\textquotedblright{}.

One of the shortest words in this book: \ta{ப} with the \textbf{\ta{ை}} sign, and that is all of it.

A \ta{பை} is what everything you \ta{வாங்கு} goes into. It is the last piece of the shop, and the one you carry home.

Five words now stand in a row you can walk through in order: the \ta{கடை} you go to, the \ta{விலை} you ask for, the \ta{பணம்} you hand over, the \ta{வாங்கு} that does it, and the \ta{பை} it all leaves in.

\begin{grammarlens}[title={What you've built}]
Five, and the shop is closed.
\end{grammarlens}

\subsection*{Your turn}
\begin{itemize}
\raggedright
  \item \emph{Say it:} \emph{pai}
  \item \emph{Say it:} all five in the order they happen — \emph{kaḍai}, \emph{vilai}, \emph{paṇam}, \emph{vāṅgu}, \emph{pai}
  \item \emph{Recall:} say \emph{paṇam}, then say \emph{vāṅgu}, then say \emph{pai}
\end{itemize}

\begin{tcolorbox}[breakable,colback=teal!4,colframe=teal!35!black,title={Before you move on}]
What does \ta{பை} mean, and which sign is on it? (\textquotedblleft{}A bag\textquotedblright{}; the \textbf{\ta{ை}}.)
\end{tcolorbox}
